\documentclass[10pt,a4paper]{amsart}
\usepackage{amsmath,amssymb,amsthm,graphicx}
\usepackage{verbatim}
\newenvironment{reference}{\comment}{\endcomment}
%\newenvironment{reference}{}{}
\newenvironment{slogan}{\comment}{\endcomment}
\newenvironment{history}{\comment}{\endcomment}

% For commutative diagrams we use Xy-pic
\usepackage[all]{xy}

% We use 2cell for 2-commutative diagrams.
\xyoption{2cell}
% Theorem environments.
%
\theoremstyle{plain}
\newtheorem{theorem}[subsection]{Theorem}
\newtheorem{proposition}[subsection]{Proposition}
\newtheorem{lemma}[subsection]{Lemma}

\theoremstyle{definition}
\newtheorem{definition}[subsection]{Definition}
\newtheorem{example}[subsection]{Example}
\newtheorem{exercise}[subsection]{Exercise}
\newtheorem{situation}[subsection]{Situation}

\theoremstyle{remark}
\newtheorem{remark}[subsection]{Remark}
\newtheorem{remarks}[subsection]{Remarks}

\numberwithin{equation}{subsection}

% Macros
%
\def\lim{\mathop{\mathrm{lim}}\nolimits}
\def\colim{\mathop{\mathrm{colim}}\nolimits}
\def\Spec{\mathop{\mathrm{Spec}}}
\def\Hom{\mathop{\mathrm{Hom}}\nolimits}
\def\Ext{\mathop{\mathrm{Ext}}\nolimits}
\def\SheafHom{\mathop{\mathcal{H}\!\mathit{om}}\nolimits}
\def\SheafExt{\mathop{\mathcal{E}\!\mathit{xt}}\nolimits}
\def\Sch{\mathit{Sch}}
\def\Mor{\mathop{\mathrm{Mor}}\nolimits}
\def\Ob{\mathop{\mathrm{Ob}}\nolimits}
\def\Sh{\mathop{\mathit{Sh}}\nolimits}
\def\NL{\mathop{N\!L}\nolimits}
\def\CH{\mathop{\mathrm{CH}}\nolimits}
\def\proetale{{pro\text{-}\acute{e}tale}}
\def\etale{{\acute{e}tale}}
\def\QCoh{\mathit{QCoh}}
\def\Ker{\mathop{\mathrm{Ker}}}
\def\Im{\mathop{\mathrm{Im}}}
\def\Coker{\mathop{\mathrm{Coker}}}
\def\Coim{\mathop{\mathrm{Coim}}}

% Boxtimes
%
\DeclareMathSymbol{\boxtimes}{\mathbin}{AMSa}{"02}

%
% Macros for moduli stacks/spaces
%
\def\QCohstack{\mathcal{QC}\!\mathit{oh}}
\def\Cohstack{\mathcal{C}\!\mathit{oh}}
\def\Spacesstack{\mathcal{S}\!\mathit{paces}}
\def\Quotfunctor{\mathrm{Quot}}
\def\Hilbfunctor{\mathrm{Hilb}}
\def\Curvesstack{\mathcal{C}\!\mathit{urves}}
\def\Polarizedstack{\mathcal{P}\!\mathit{olarized}}
\def\Complexesstack{\mathcal{C}\!\mathit{omplexes}}
% \Pic is the operator that assigns to X its picard group, usage \Pic(X)
% \Picardstack_{X/B} denotes the Picard stack of X over B
% \Picardfunctor_{X/B} denotes the Picard functor of X over B
\def\Pic{\mathop{\mathrm{Pic}}\nolimits}
\def\Picardstack{\mathcal{P}\!\mathit{ic}}
\def\Picardfunctor{\mathrm{Pic}}
\def\Deformationcategory{\mathcal{D}\!\mathit{ef}}
\usepackage{cleveref}
\crefname{originalreference}{original reference}{original references}
\Crefname{originalreference}{Original reference}{Original references}
\makeatletter
\newlabel{situation-setup}{{02QL}{source}{}{sourcepointer-1}{}}
\newlabel{situation-setup@cref}{{[originalreference][0][]02QL}{[1][0][]source}{}{}{}}
\makeatother
\crefname{originallemma}{original lemma}{original lemmas}
\Crefname{originallemma}{Original lemma}{Original lemmas}
\makeatletter
\newlabel{more-algebra-lemma-blowup-regular-sequence}{{0BIQ}{source}{}{sourcepointer-2}{}}
\newlabel{more-algebra-lemma-blowup-regular-sequence@cref}{{[originallemma][0][]0BIQ}{[1][0][]source}{}{}{}}
\makeatother
\crefname{originallemma}{original lemma}{original lemmas}
\Crefname{originallemma}{Original lemma}{Original lemmas}
\makeatletter
\newlabel{divisors-lemma-quasi-regular-immersion}{{063M}{source}{}{sourcepointer-3}{}}
\newlabel{divisors-lemma-quasi-regular-immersion@cref}{{[originallemma][0][]063M}{[1][0][]source}{}{}{}}
\makeatother
\crefname{originallemma}{original lemma}{original lemmas}
\Crefname{originallemma}{Original lemma}{Original lemmas}
\makeatletter
\newlabel{lemma-compare-gysin-base-change}{{0FV9}{source}{}{sourcepointer-4}{}}
\newlabel{lemma-compare-gysin-base-change@cref}{{[originallemma][0][]0FV9}{[1][0][]source}{}{}{}}
\makeatother
\crefname{originallemma}{original lemma}{original lemmas}
\Crefname{originallemma}{Original lemma}{Original lemmas}
\makeatletter
\newlabel{lemma-lci-gysin-flat}{{0FF4}{source}{}{sourcepointer-5}{}}
\newlabel{lemma-lci-gysin-flat@cref}{{[originallemma][0][]0FF4}{[1][0][]source}{}{}{}}
\makeatother
\crefname{originallemma}{original lemma}{original lemmas}
\Crefname{originallemma}{Original lemma}{Original lemmas}
\makeatletter
\newlabel{more-morphisms-lemma-NL-smooth}{{0D0N}{source}{}{sourcepointer-6}{}}
\newlabel{more-morphisms-lemma-NL-smooth@cref}{{[originallemma][0][]0D0N}{[1][0][]source}{}{}{}}
\makeatother
\crefname{originallemma}{original lemma}{original lemmas}
\Crefname{originallemma}{Original lemma}{Original lemmas}
\makeatletter
\newlabel{more-morphisms-lemma-get-triangle-NL}{{0FV3}{source}{}{sourcepointer-7}{}}
\newlabel{more-morphisms-lemma-get-triangle-NL@cref}{{[originallemma][0][]0FV3}{[1][0][]source}{}{}{}}
\makeatother
\crefname{originallemma}{original lemma}{original lemmas}
\Crefname{originallemma}{Original lemma}{Original lemmas}
\makeatletter
\newlabel{more-morphisms-lemma-NL-immersion}{{0FV2}{source}{}{sourcepointer-8}{}}
\newlabel{more-morphisms-lemma-NL-immersion@cref}{{[originallemma][0][]0FV2}{[1][0][]source}{}{}{}}
\makeatother
\title{The blow up formula: the restriction of the Gysin map of a blowup in a regular center is the top Chern class of the kernel of the conormal comparison}
\author{The Stacks Project Authors}
\date{stacks-project:chow.tex:local-snapshot-20261012; GNU FDL 1.2 or later}
\begin{document}
\maketitle
\noindent\textit{Editorial note.} All delivered mathematical text is the author's own, copied from the pinned Stacks Project chapter file chow.tex; no statement or proof step was rewritten, reordered or removed.\par\medskip
\noindent\textit{Editorial note.} Preamble provenance: the pinned chapter file opens with an input of the shared preamble and carries no preamble text of its own, so the project's single Stacks-wide preamble staged at sources/stacks-chow/preamble.tex is excerpted verbatim as two spans, lines 14-18 (Xy-pic setup) and lines 154-217 (the native theorem declarations and mathematical macros). Nothing was added, removed or redefined; the wrapper class and the amsmath/amssymb/amsthm/graphicx base packages are supplied by the assembler.\par\medskip
\noindent\textit{Editorial note.} The counted statement and proof contain no cross reference, no citation and no included file, so the delivered body is self contained.\par\medskip
\noindent\textit{Editorial note.} Canonical location of the counted proposition: Stacks tag 0FVA (sources/stacks-chow/pinned-original-source.tex line 13808).\par\medskip
\noindent\textit{Allowed standard prerequisites (editorial).}\begin{itemize}
\item The blowup in a regular sequence is a local complete intersection morphism. More on Algebra, Lemma labelled more-algebra-lemma-blowup-regular-sequence, cited at proof lines 13827-13829. Makes the Gysin map b\textasciicircum{}! available and the statement meaningful.
\item The conormal algebra of a regular immersion is the symmetric algebra. Divisors, Lemma labelled divisors-lemma-quasi-regular-immersion, cited at proof lines 13832-13837. Identifies E with the projective space bundle of the conormal sheaf, hence E -> Z smooth.
\item Comparison of Gysin maps under a base change, and equality of the Gysin and the flat pullback for local complete intersection morphisms. Same chapter, Lemma labelled lemma-compare-gysin-base-change, cited at proof lines 13842-13846, and Lemma labelled lemma-lci-gysin-flat, cited at proof lines 13846-13848. Produces res(b\textasciicircum{}!) = c\_top(C\textasciicircum{}vee) o pi\textasciicircum{}! and turns pi\textasciicircum{}! into pi\textasciicircum{}*.
\item The exact sequences of the normal cone for a closed immersion and a smooth morphism. More on Morphisms, Lemmas labelled more-morphisms-lemma-NL-smooth, more-morphisms-lemma-get-triangle-NL and more-morphisms-lemma-NL-immersion, cited at proof lines 13852-13862. Identifies the two kernels F and C.
\end{itemize}\medskip
\ifdefined\hypertarget\hypertarget{sourcepointer-1}{}\fi
\noindent\textit{Source reference.} Original reference 02QL: sources/stacks-chow/pinned-original-source.tex, source lines 1372--1372. Same-chapter pointer into the author environment situation at source line 1372; invoked at proof line None of the counted proof. It is retained only as a navigation pointer; every load-bearing step is inside the delivered span.\par\medskip
\ifdefined\hypertarget\hypertarget{sourcepointer-2}{}\fi
\noindent\textit{Source reference.} Original lemma 0BIQ: sources/stacks-chow/tags.txt, source lines 14498--14498. Cross-chapter standard prerequisite located in the project tag list staged in the local snapshot; invoked at proof line 13829 of the counted proof. It is retained only as a navigation pointer; every load-bearing step is inside the delivered span.\par\medskip
\ifdefined\hypertarget\hypertarget{sourcepointer-3}{}\fi
\noindent\textit{Source reference.} Original lemma 063M: sources/stacks-chow/tags.txt, source lines 7836--7836. Cross-chapter standard prerequisite located in the project tag list staged in the local snapshot; invoked at proof line 13835 of the counted proof. It is retained only as a navigation pointer; every load-bearing step is inside the delivered span.\par\medskip
\ifdefined\hypertarget\hypertarget{sourcepointer-4}{}\fi
\noindent\textit{Source reference.} Original lemma 0FV9: sources/stacks-chow/pinned-original-source.tex, source lines 13738--13738. Same-chapter pointer into the author environment lemma at source line 13738; invoked at proof line 13840 of the counted proof. It is retained only as a navigation pointer; every load-bearing step is inside the delivered span.\par\medskip
\ifdefined\hypertarget\hypertarget{sourcepointer-5}{}\fi
\noindent\textit{Source reference.} Original lemma 0FF4: sources/stacks-chow/pinned-original-source.tex, source lines 13547--13547. Same-chapter pointer into the author environment lemma at source line 13547; invoked at proof line 13846 of the counted proof. It is retained only as a navigation pointer; every load-bearing step is inside the delivered span.\par\medskip
\ifdefined\hypertarget\hypertarget{sourcepointer-6}{}\fi
\noindent\textit{Source reference.} Original lemma 0D0N: sources/stacks-chow/tags.txt, source lines 16318--16318. Cross-chapter standard prerequisite located in the project tag list staged in the local snapshot; invoked at proof line 13853 of the counted proof. It is retained only as a navigation pointer; every load-bearing step is inside the delivered span.\par\medskip
\ifdefined\hypertarget\hypertarget{sourcepointer-7}{}\fi
\noindent\textit{Source reference.} Original lemma 0FV3: sources/stacks-chow/tags.txt, source lines 19798--19798. Cross-chapter standard prerequisite located in the project tag list staged in the local snapshot; invoked at proof line 13856 of the counted proof. It is retained only as a navigation pointer; every load-bearing step is inside the delivered span.\par\medskip
\ifdefined\hypertarget\hypertarget{sourcepointer-8}{}\fi
\noindent\textit{Source reference.} Original lemma 0FV2: sources/stacks-chow/tags.txt, source lines 19797--19797. Cross-chapter standard prerequisite located in the project tag list staged in the local snapshot; invoked at proof line 13857 of the counted proof. It is retained only as a navigation pointer; every load-bearing step is inside the delivered span.\par\medskip
\begin{lemma}[Blow up formula]
\label{lemma-blow-up-formula}
Let $(S, \delta)$ be as in Situation \ref{situation-setup}.
Let $i : Z \to X$ be a regular closed immersion of schemes
locally of finite type over $S$. Let $b : X' \to X$ be the
blowing up with center $Z$. Picture
$$
\xymatrix{
E \ar[r]_j \ar[d]_\pi & X' \ar[d]^b \\
Z \ar[r]^i & X
}
$$
Assume that the Gysin map exists for $b$. Then we have
$$
res(b^!) = c_{top}(\mathcal{F}^\vee) \circ \pi^*
$$
in $A^*(E \to Z)$ where $\mathcal{F}$ is the kernel of the canonical map
$\pi^*\mathcal{C}_{Z/X} \to \mathcal{C}_{E/X'}$.
\end{lemma}

\begin{proof}
Observe that the morphism $b$ is a local complete intersection morphism
by More on Algebra, Lemma \ref{more-algebra-lemma-blowup-regular-sequence}
and hence the statement makes sense. Since $Z \to X$ is a regular
immersion (and hence a fortiori quasi-regular) we see that $\mathcal{C}_{Z/X}$
is finite locally free and the map
$\text{Sym}^*(\mathcal{C}_{Z/X}) \to \mathcal{C}_{Z/X, *}$
is an isomorphism, see
Divisors, Lemma \ref{divisors-lemma-quasi-regular-immersion}.
Since $E = \text{Proj}(\mathcal{C}_{Z/X, *})$ we conclude
that $E = \mathbf{P}(\mathcal{C}_{Z/X})$
is a projective space bundle over $Z$.
Thus $E \to Z$ is smooth and certainly a local complete intersection
morphism. Thus Lemma \ref{lemma-compare-gysin-base-change}
applies and we see that
$$
res(b^!) = c_{top}(\mathcal{C}^\vee) \circ \pi^!
$$
with $\mathcal{C}$ as in the statement there.
Of course $\pi^* = \pi^!$ by Lemma \ref{lemma-lci-gysin-flat}.
It remains to show that $\mathcal{F}$ is equal to
the kernel $\mathcal{C}$ of the map
$H^{-1}(j^*\NL_{X'/X}) \to H^{-1}(\NL_{E/Z})$.

\medskip\noindent
Since $E \to Z$ is smooth we have $H^{-1}(\NL_{E/Z}) = 0$, see
More on Morphisms, Lemma \ref{more-morphisms-lemma-NL-smooth}.
Hence it suffices to show that $\mathcal{F}$ can be identified
with $H^{-1}(j^*\NL_{X'/X})$. By More on Morphisms, Lemmas
\ref{more-morphisms-lemma-get-triangle-NL} and
\ref{more-morphisms-lemma-NL-immersion} we have an exact sequence
$$
0 \to H^{-1}(j^*\NL_{X'/X}) \to H^{-1}(\NL_{E/X}) \to
\mathcal{C}_{E/X'} \to \ldots
$$
By the same lemmas applied to $E \to Z \to X$ we obtain an isomorphism
$\pi^*\mathcal{C}_{Z/X} = H^{-1}(\pi^*\NL_{Z/X}) \to H^{-1}(\NL_{E/X})$.
Thus we conclude.
\end{proof}

\end{document}
